% Plan de trabajo y ataque — Hexápodo (versión paralela de la Evaluación 1)
% Compilar con XeLaTeX (dos pasadas):  xelatex plan_ataque.tex
\documentclass[11pt,a4paper]{article}

\usepackage[spanish,es-tabla]{babel}
\usepackage{fontspec}
\setmainfont{Lato}
\usepackage{unicode-math}
\setmathfont{TeX Gyre DejaVu Math}
\usepackage[a4paper,margin=2.2cm,top=2.6cm,bottom=2.4cm,headheight=28pt]{geometry}
\usepackage{xcolor,graphicx,booktabs,tabularx,array,enumitem,amsmath,fancyhdr,lastpage,titlesec,multirow,colortbl}
\usepackage[most]{tcolorbox}
\usepackage{tikz}
\usetikzlibrary{positioning,arrows.meta,calc,fit,shapes.geometric}
\usepackage[hidelinks]{hyperref}
\usepackage{caption}
\usepackage{float}
\usepackage{needspace}
\setcounter{tocdepth}{2}

% ---------------- Paleta ----------------
\definecolor{navy}{HTML}{002F6C}      % USACH 294
\definecolor{naranjo}{HTML}{EA7600}   % USACH 716
\definecolor{gris}{HTML}{5F6B7A}
\definecolor{fondo}{HTML}{FFFFFF}
% Color = tipo de trabajo
\definecolor{cModelo}{HTML}{1F5FAD}
\definecolor{cControl}{HTML}{4B3B9A}
\definecolor{cSim}{HTML}{2E7D32}
\definecolor{cSens}{HTML}{00838F}
\definecolor{cMicro}{HTML}{EA7600}
\definecolor{cInteg}{HTML}{A6790A}
\definecolor{cEstudio}{HTML}{7B4A2A}
\definecolor{cEntrega}{HTML}{5F6B7A}
% Semáforo
\definecolor{sRojo}{HTML}{C62828}
\definecolor{sAmbar}{HTML}{E69500}
\definecolor{sVerde}{HTML}{43A047}

\captionsetup{font=small,labelfont={bf,color=navy}}
\setlength{\parindent}{0pt}
\setlength{\parskip}{0.45em}
\renewcommand{\arraystretch}{1.25}
\setlist{itemsep=0.15em,topsep=0.3em}

% Etiqueta de color (pill)
\newcommand{\pill}[2]{{\setlength{\fboxsep}{1.6pt}\colorbox{#1}{\textcolor{white}{\footnotesize\bfseries\,#2\,}}}}
\newcommand{\Ez}{\pill{cModelo}{E0}}
\newcommand{\Eu}{\pill{cControl}{E1}}
\newcommand{\Ed}{\pill{cModelo}{E2}}
\newcommand{\Et}{\pill{cSim}{E3}}
\newcommand{\Ec}{\pill{cControl}{E4}}
\newcommand{\Eci}{\pill{cSens}{E5}}
\newcommand{\Es}{\pill{cMicro}{E6}}
\newcommand{\Esi}{\pill{cInteg}{E7}}
\newcommand{\Eo}{\pill{cEstudio}{E8}}
\newcommand{\En}{\pill{cEntrega}{E9}}

% Recuadro destacado: mismo fondo de la hoja, barra lateral y título en el color de acento
\newtcolorbox{destacado}[2][naranjo]{enhanced,breakable,colback=fondo,colframe=fondo,boxrule=0pt,
  borderline west={3pt}{0pt}{#1},left=8pt,right=4pt,top=2pt,bottom=2pt,
  title={#2},coltitle=#1,colbacktitle=fondo,fonttitle=\bfseries,attach title to upper={\par\smallskip}}

% Títulos
\titleformat{\section}{\color{navy}\Large\bfseries}{\thesection}{0.6em}{}[\vspace{-0.3em}{\color{naranjo}\rule{\linewidth}{1.2pt}}]
\titleformat{\subsection}{\color{navy}\large\bfseries}{\thesubsection}{0.5em}{}
\titleformat{\subsubsection}{\color{gris}\normalsize\bfseries}{}{0em}{}
\titlespacing*{\section}{0pt}{1.2em}{0.6em}

% Encabezado y pie
\pagestyle{fancy}
\fancyhf{}
\fancyhead[L]{\color{navy}\small\bfseries Hexápodo --- Plan de trabajo y ataque}
\fancyhead[R]{\includegraphics[height=20pt]{logo_small.png}}
\fancyfoot[L]{\color{gris}\footnotesize Open Robotics · USACH · 2.º semestre 2026}
\fancyfoot[R]{\color{gris}\footnotesize Página \thepage\ de \pageref{LastPage}}
\renewcommand{\headrulewidth}{0.8pt}
\renewcommand{\headrule}{{\color{navy}\hrule height \headrulewidth}}

% Ficha de etapa
\newcommand{\ficha}[4]{% color, tipo, cuándo, pauta
  \begin{tabularx}{\linewidth}{@{}l Y@{}}
  \textbf{Tipo:} & \pill{#1}{#2}\\
  \textbf{Cuándo:} & #3\\
  \textbf{Puntos de la pauta:} & #4\\
  \textbf{Responsable:} & César (revisa, estudia y decide) · Apoyo: Claude (código)\\
  \end{tabularx}\par}

\newcolumntype{Y}{>{\raggedright\arraybackslash}X}

\begin{document}

% ====================== PORTADA ======================
\begin{titlepage}
\newgeometry{margin=0pt}
\begin{tikzpicture}[remember picture,overlay]
  \fill[navy] (current page.north west) rectangle ([yshift=-11cm]current page.north east);
  \fill[naranjo] ([yshift=-11cm]current page.north west) rectangle ([yshift=-11.35cm]current page.north east);
  \node[anchor=north west] at ([xshift=2cm,yshift=-1.6cm]current page.north west) {\colorbox{white}{\includegraphics[height=2.6cm]{usach_logo.jpg}}};
  \node[anchor=west,text=white,font=\fontsize{30}{34}\selectfont\bfseries,align=left] at ([xshift=2cm,yshift=-6.4cm]current page.north west) {Hexápodo con ROS\,2};
  \node[anchor=west,text=white,font=\Large,align=left] at ([xshift=2cm,yshift=-7.9cm]current page.north west) {Plan de trabajo y ataque del proyecto};
  \node[anchor=west,text=naranjo!40,font=\large,align=left] at ([xshift=2cm,yshift=-9cm]current page.north west) {Versión paralela de la Evaluación 1 · Estudio de dimensiones};
  \node[anchor=north west,font=\large,align=left,text=navy] at ([xshift=2cm,yshift=-13cm]current page.north west) {%
    \begin{tabular}{@{}ll@{}}
    \textbf{Asignatura} & Open Robotics --- 2.º semestre 2026\\[0.3em]
    \textbf{Estudiante} & Cesar Olivares Lavalovich\\[0.3em]
    \textbf{Profesor} & José Pascal\\[0.3em]
    \textbf{Fecha} & 26 de noviembre de 2026\\[0.3em]
    \textbf{Versión} & v1\\
    \end{tabular}};
  \node[anchor=south west,font=\small,text=gris,align=left,text width=17cm] at ([xshift=2cm,yshift=2cm]current page.south west) {%
    Universidad de Santiago de Chile. Documento de trabajo personal: guía para ejecutar, estudiar y presentar el proyecto.
    La entrega oficial de la Evaluación 1 es el brazo Octavius; este proyecto cumple la misma pauta con un hexápodo.};
\end{tikzpicture}
\restoregeometry
\end{titlepage}

\tableofcontents
\clearpage

% ====================== 1. META ======================
\Needspace{14\baselineskip}
\section{Meta y alcance}

\begin{destacado}{Meta}
De aquí al \textbf{domingo 29 de noviembre de 2026} (\emph{fecha supuesta, por confirmar}), tener un hexápodo simulado en ROS\,2 Jazzy que cumpla \textbf{todos los puntos de la pauta}. Además, tener las \textbf{cotas ideales} de patas y cuerpo, calculadas a partir de los componentes reales, para diseñar y maquinar después las piezas.
\end{destacado}

\textbf{Dentro del alcance}
\begin{itemize}
  \item Todo lo que pide la pauta, en simulación:
  \begin{itemize}
    \item workspace y paquetes;
    \item URDF/XACRO;
    \item RViz y Gazebo;
    \item dos sensores;
    \item \texttt{ros2\_control} con teclado;
    \item micro-ROS;
    \item launch;
    \item presentación, tutorial y \texttt{.zip}.
  \end{itemize}
  \item Modelo con las masas y los datos de hoja de los componentes reales: 18 MG996R, Raspberry Pi 4, 2 PCA9685, LiPo 3S 5200 mAh, IMU y cámara.
  \item Estudio de dimensiones: largos de coxa, fémur y tibia, radio del hexágono y layout de los componentes.
\end{itemize}

\textbf{Fuera del alcance}
\begin{itemize}
  \item Construir, cablear o poner en marcha el robot real.
  \item Las piezas CAD hechas a ojo (\texttt{base\_hexapodo}, \texttt{BICEP}, \texttt{PATA\_1}, \texttt{Tapa superior}). Se descartan y se rediseñan con los resultados del estudio.
\end{itemize}

\begin{destacado}[sRojo]{Aviso: punto 9 de la pauta (micro-ROS)}
La pauta pide conectar una ESP32 por micro-ROS y mover articulaciones con ella. La ESP32 todavía no está comprada. Esta semana se deja listo todo lo demás:
\begin{itemize}
  \item el paquete;
  \item el tópico;
  \item el nodo que convierte el joystick en movimiento;
  \item el firmware compilado.
\end{itemize}
Sin la ESP32, el robot se mueve publicando el tópico a mano. \textbf{El punto solo queda cerrado al conectar la placa real.} Si no se compra antes de presentar, hay que decirlo en la presentación.
\end{destacado}

% ====================== 2. OBJETIVOS ======================
\Needspace{14\baselineskip}
\section{Objetivos}

\begin{table}[H]
\centering
\begin{tabularx}{\linewidth}{@{}l l Y@{}}
\toprule
\textbf{Objetivo} & \textbf{Fecha objetivo} & \textbf{Descripción del objetivo}\\
\midrule
\pill{cModelo}{O1} Modelo en RViz & \textbf{vie 27 nov} & URDF/XACRO del hexápodo (18 articulaciones, componentes, sensores) visible en RViz con 18 sliders, RobotModel y TF. Pauta 1--5.\\
\pill{cSim}{O2} Caminar en Gazebo & sáb 28 nov & Se sostiene de pie en Gazebo y camina con el teclado mediante \texttt{ros2\_control}. Pauta 6 y 8.\\
\pill{cSens}{O3} Sensores, launch y micro-ROS & sáb 28 nov & IMU y cámara publicando y visibles en RViz. Un launch levanta todo. Paquete micro-ROS listo para la ESP32. Pauta 7, 9 (parcial), 10 y 11.\\
\pill{cEstudio}{O4} Cotas ideales & dom 29 nov & Largos de patas, radio del cuerpo y layout que cumplen torque, alcance, estabilidad y sin choques. Hoja de especificación para maquinar.\\
\pill{cEntrega}{O5} Entrega & dom 29 nov & Presentación (PDF + editable), tutorial (PDF + editable) y \texttt{.zip} con el nombre pedido. Pauta 12--14.\\
\bottomrule
\end{tabularx}
\caption{Objetivos de la semana. Intensidad del color: lo que cierra antes va más intenso en el cronograma (sección~\ref{sec:cronograma}).}
\end{table}

% ====================== 3. PAUTA ======================
\Needspace{14\baselineskip}
\section{Trazabilidad con la pauta}

La pauta es la regla del proyecto: ningún cambio puede dejar uno de estos puntos sin cumplir. Los puntos se numeran P1--P14.

\begin{table}[H]
\centering\small
\begin{tabularx}{\linewidth}{@{}l Y l Y@{}}
\toprule
\textbf{Punto} & \textbf{Qué pide} & \textbf{Etapa} & \textbf{Evidencia}\\
\midrule
P1 & Paquete con sus directorios & \Ez & \texttt{colcon build} sin errores\\
P2 & Robot articulado en URDF/XACRO & \Ed & \texttt{check\_urdf} OK, árbol con \texttt{urdf\_to\_graphviz}\\
P3 & Modelo en RViz & \Ed & Captura de RViz\\
P4 & \texttt{joint\_state\_publisher} & \Ed & 18 sliders mueven las patas\\
P5 & Percepción en RViz: RobotModel y TF & \Ed\ \Eci & Captura con RobotModel, TF, imagen e IMU\\
P6 & Modelo en Gazebo & \Et & De pie $\geq$ 60 s sin hundirse\\
P7 & $\geq$ 2 plugins de sensores & \Eci & \texttt{ros2 topic hz} de IMU y cámara\\
P8 & \texttt{ros2\_control} + teclado & \Et\ \Ec & Camina y gira con teclas\\
P9 & micro-ROS & \Es & Tópico a mano esta semana; ESP32 real al comprarla\\
P10 & Launchfile & \Esi & Un comando levanta todo\\
P11 & Estructura de carpetas & \Ez\ \Esi & Revisión del árbol del repositorio\\
P12 & Presentación PDF + editable, $\leq$ 15 min & \En & Checklist de formato (sección~\ref{sec:entrega})\\
P13 & Tutorial PDF + editable & \En & \texttt{DOCUMENTACION.md} $\rightarrow$ \texttt{.docx} + \texttt{.pdf}\\
P14 & \texttt{.zip} \texttt{nombre1\_apellido1\_año\_semestre} & \En & Descomprimir en carpeta limpia y compilar\\
\bottomrule
\end{tabularx}
\caption{Puntos de la pauta, etapa que los cumple y cómo se comprueba.}
\end{table}

% ====================== 4. DATOS ======================
\Needspace{14\baselineskip}
\section{Datos del robot}

\subsection{Componentes y masas}

Lo medido o sacado de una hoja de datos es un \textbf{hecho}. Lo marcado como \emph{hipótesis} es una estimación propia, y el estudio de dimensiones la revisa.

\begin{table}[H]
\centering\small
\begin{tabularx}{\linewidth}{@{}l r Y l@{}}
\toprule
\textbf{Componente} & \textbf{Masa} & \textbf{Datos usados} & \textbf{Origen}\\
\midrule
18 × MG996R (eje coaxial) & 18 × 55 g & Cuerpo 40,7 × 19,7 × 42,9 mm. Bloqueo 1,08 N·m (6 V). 0,14 s/60° (6 V) & Pesado · hoja\\
Adaptadores coaxiales A/B & --- & A: 49 × 69 × 20 mm; B: 40 × 49 × 20 mm & STEP del repositorio\\
LiPo 3S 5200 mAh 60C & 271 g & 11,1 V nominal; medidas por confirmar & Pesada\\
Raspberry Pi 4 & 86 g & 85 × 56 mm & Pesada · hoja\\
2 × PCA9685 & 2 × 9 g & 62 × 25 mm; 16 canales c/u & Pesado · STEP\\
IMU BNO085 (por comprar) & 3 g & Orientación + aceleración + giro & \emph{Hipótesis}\\
Cámara Pi Camera Module 3 & 4 g & IMX708, campo de visión horizontal 66° & Hoja\\
Estructura impresa & 400 g & Placas del cuerpo y piezas de patas & \emph{Hipótesis}\\
Cables y reguladores & 100 g & --- & \emph{Hipótesis}\\
\midrule
\textbf{Total} & \textbf{$\approx$ 1,87 kg} & Peso $W \approx 18{,}4$ N & \\
\bottomrule
\end{tabularx}
\caption{Componentes del robot y datos que entran al modelo.}
\end{table}

\subsection{Primera cota: cuánto puede alejarse el pie}

Con la marcha trípode, 3 patas sostienen todo el peso. Cada una recibe una fuerza vertical:
\begin{equation}
F = \frac{m\,g}{3} = \frac{1{,}87 \cdot 9{,}81}{3} \approx 6{,}1\ \text{N}
\end{equation}
El servo más cargado es el del fémur. Su torque es esa fuerza por la distancia horizontal $d$ entre su eje y el pie. Para dejar margen, se le exige como máximo la mitad del torque de bloqueo (factor de seguridad 2):
\begin{equation}
F\,d \;\leq\; \tau_d = \frac{\tau_{\text{bloqueo}}}{2} = \frac{1{,}08}{2} = 0{,}54\ \text{N·m}
\quad\Rightarrow\quad d \leq 8{,}8\ \text{cm}\ \ (7{,}5\ \text{cm a 4,8 V})
\end{equation}

\begin{destacado}{Por qué importa}
Con estos servos, una pata larga y muy abierta no es viable. Esta cota sirve de punto de partida para \texttt{hexapod\_params.xacro} desde el primer día. El estudio de dimensiones (\Eo) la refina con el peso propio de las patas y con la dinámica medida en Gazebo.
\end{destacado}

% ====================== 5. ARQUITECTURA ======================
\Needspace{14\baselineskip}
\section{Arquitectura del sistema}

\subsection{Cuerpo y patas}

El cuerpo es un \textbf{hexágono regular horizontal}, con una pata en cada vértice. El robot avanza hacia una \textbf{cara} y no hacia un vértice. Así quedan 3 patas simétricas por lado, a $\pm30^\circ$, $\pm90^\circ$ y $\pm150^\circ$ del eje de avance, y la cámara queda centrada en la cara frontal. El radio $R$ del hexágono es un resultado del estudio.

\begin{figure}[H]
\centering
\begin{tikzpicture}[scale=1.0]
  \def\R{2.2}
  \draw[navy,line width=1.2pt,fill=navy!5] (30:\R) -- (90:\R) -- (150:\R) -- (210:\R) -- (270:\R) -- (330:\R) -- cycle;
  \foreach \a in {30,90,150,210,270,330}{
    \draw[cControl,line width=2pt] (\a:\R) -- (\a:{\R+1.3});
    \fill[cControl] (\a:{\R+1.3}) circle (2.5pt);
  }
  \draw[-{Stealth[length=3mm]},naranjo,line width=1.2pt] (0,0) -- (0:1.6) node[right,font=\small\bfseries]{avance $x$};
  \fill[cSens] (0:{\R*0.866}) circle (3pt);
  \node[font=\scriptsize,cSens,right] at ($(0:{\R*0.866})+(0.05,-0.3)$) {cámara};
  \fill[cSens] (0,0) circle (2.5pt) node[below left,font=\scriptsize]{IMU};
  \draw[dashed,gris] (0,0) -- (30:\R) node[midway,above left,font=\scriptsize]{$R$};
  \draw[-{Stealth[length=3mm]},gris,line width=0.8pt] (0,0) -- (90:0.9) node[left,font=\scriptsize]{$y$};
  \node[font=\scriptsize\bfseries] at (30:{\R+1.75}) {LF};
  \node[font=\scriptsize\bfseries] at (90:{\R+1.75}) {LM};
  \node[font=\scriptsize\bfseries] at (150:{\R+1.75}) {LR};
  \node[font=\scriptsize\bfseries] at (330:{\R+1.75}) {RF};
  \node[font=\scriptsize\bfseries] at (270:{\R+1.75}) {RM};
  \node[font=\scriptsize\bfseries] at (210:{\R+1.75}) {RR};
\end{tikzpicture}
\caption{Cuerpo hexagonal regular visto desde arriba (convención REP-103: $x$ hacia el frente, $y$ hacia la izquierda). Las patas salen de los vértices y el frente es una cara.}
\end{figure}

\subsection{Flujo de datos}

\begin{figure}[H]
\centering
\begin{tikzpicture}[
  n/.style={draw=navy,fill=navy!6,rounded corners=2pt,font=\scriptsize\bfseries,align=center,minimum height=7mm,inner sep=3pt},
  t/.style={draw=gris,fill=gris!10,font=\tiny\ttfamily,align=center,minimum height=6mm,inner sep=2pt},
  m/.style={draw=cMicro,fill=cMicro!10,rounded corners=2pt,font=\scriptsize\bfseries,align=center,minimum height=7mm,inner sep=3pt},
  a/.style={-{Stealth[length=2mm]},line width=0.7pt,draw=black!70},
  node distance=4mm and 5mm]
  \node[m] (esp) at (0,-1.8) {ESP32 +\\joystick};
  \node[m] (ag) at (2.7,-1.8) {micro\_ros\\agent};
  \node[t] (joy) at (5.3,-1.8) {/hexapod/\\joystick};
  \node[n] (j2t) at (7.9,-1.8) {joy\_to\\twist};
  \node[n] (kb) at (5.3,0) {keyboard\\teleop};
  \node[t] (cmd) at (7.9,0) {/hexapod/\\cmd\_vel};
  \node[n] (gait) at (10.6,0) {gait\_node\\(marcha + IK)};
  \node[t] (lc) at (10.6,-1.8) {/legs\_controller/\\commands};
  \node[n] (gz) at (13.6,-1.8) {gz\_ros2\_control\\(Gazebo)};
  \node[t] (js) at (13.6,0) {/joint\_states};
  \node[n] (rviz) at (13.6,1.6) {RViz};
  \draw[a] (esp) -- node[below,font=\tiny]{USB} (ag);
  \draw[a] (ag) -- (joy);
  \draw[a] (joy) -- (j2t);
  \draw[a] (j2t) -- (cmd);
  \draw[a] (kb) -- (cmd);
  \draw[a] (cmd) -- (gait);
  \draw[a] (gait) -- (lc);
  \draw[a] (lc) -- (gz);
  \draw[a] (gz) -- (js);
  \draw[a] (js) -- (rviz);
\end{tikzpicture}
\caption{Del operador a las articulaciones. En naranjo, lo que depende de la ESP32.}
\end{figure}

Hay un solo nodo, \texttt{gait\_node}, que calcula los 18 ángulos. El teclado y el joystick solo piden una \textbf{velocidad del cuerpo}, un mensaje \texttt{geometry\_msgs/Twist} con avance $v_x$, desplazamiento lateral $v_y$ y giro $\omega_z$. Así ambas fuentes son intercambiables y la marcha no sabe de dónde viene la orden.

\subsection{Controladores}

\begin{itemize}
  \item \texttt{joint\_state\_broadcaster}: publica \texttt{/joint\_states} con el ángulo medido de las 18 articulaciones.
  \item \texttt{legs\_controller} (\texttt{JointGroupPositionController}): recibe un vector de 18 ángulos y lo copia tal cual a las articulaciones. No interpola. La suavidad la pone la marcha, que publica a 50 Hz.
\end{itemize}

% ====================== 6. CRONOGRAMA ======================
\Needspace{14\baselineskip}
\section{Cronograma}\label{sec:cronograma}

Cada día se divide en tres bloques: mañana, tarde y noche. Hoy jueves solo queda la noche. El color indica el tipo de trabajo; la intensidad, qué tan pronto parte.

\begin{figure}[H]
\centering
\begin{tikzpicture}[x=1.12cm,y=0.55cm]
  % columnas: 1 jue-n | 2 vie-m 3 vie-t 4 vie-n | 5 sab-m 6 sab-t 7 sab-n | 8 dom-m 9 dom-t 10 dom-n
  \foreach \d/\x/\w in {Jue 26/0/1,Vie 27/1/3,Sáb 28/4/3,Dom 29/7/3}{
    \draw[gris!50] (\x,0.6) rectangle ++(\w,-11.2);
    \node[font=\small\bfseries,navy] at (\x+\w/2,1.1) {\d};
  }
  \foreach \i/\l in {0/N,1/M,2/T,3/N,4/M,5/T,6/N,7/M,8/T,9/N}{\node[font=\tiny,gris] at (\i+0.5,0.25) {\l};}
  \newcommand{\barra}[5]{% fila, inicio, fin, color, texto
    \fill[#4,rounded corners=1.5pt] (#2+0.05,-#1+0.05-0.45) rectangle (#3-0.05,-#1-0.45-0.35);
    \node[anchor=east,font=\scriptsize] at (-0.1,-#1-0.6) {#5};}
  \barra{0}{0}{1}{cModelo}{E0 Workspace y parámetros}
  \barra{1}{0}{2}{cControl}{E1 Cinemática de una pata}
  \barra{2}{1}{3}{cModelo}{E2 Modelo completo en RViz}
  \barra{3}{2}{4}{cSim}{E3 De pie en Gazebo}
  \barra{4}{4}{6}{cControl!75}{E4 Marcha y teclado}
  \barra{5}{5}{6}{cSens!75}{E5 Sensores}
  \barra{6}{6}{7}{cMicro!75}{E6 micro-ROS}
  \barra{7}{6}{7}{cInteg!75}{E7 Integración y pruebas}
  \barra{8}{7}{9}{cEstudio!55}{E8 Estudio de dimensiones}
  \barra{9}{8}{10}{cEntrega!55}{E9 Documentación y entrega}
  % línea de corte
  \draw[sRojo,dashed,line width=1pt] (8,0.6) -- (8,-11);
  \node[sRojo,font=\tiny\bfseries,anchor=south,align=center] at (8,-11.9) {línea de corte};
\end{tikzpicture}
\caption{Cronograma por bloques (N: noche, M: mañana, T: tarde). A la derecha de la línea de corte queda lo que se recorta primero si falta tiempo.}
\end{figure}

\begin{destacado}[sRojo]{Si falta tiempo: qué se recorta y en qué orden}
\begin{enumerate}
  \item La verificación de torques en Gazebo (\Eo, T4). El estudio queda con el modelo analítico.
  \item La prueba en limpio dentro de un contenedor (\Esi, T3).
  \item El desplazamiento lateral y el giro en el lugar. Se conservan avanzar, retroceder y girar avanzando.
\end{enumerate}
\textbf{Ningún punto de la pauta se recorta.} \En\ no se recorta: la documentación se escribe al cerrar cada etapa, no al final.
\end{destacado}

% ====================== 7. ETAPAS ======================
\Needspace{14\baselineskip}
\section{Etapas}

Cada etapa tiene esta estructura:
\begin{itemize}
  \item \textbf{qué hay que entender}, para estudiar y explicar en la presentación;
  \item \textbf{tareas};
  \item \textbf{entregables con su evidencia};
  \item \textbf{qué mostrar en la presentación}.
\end{itemize}
Una etapa cierra cuando todas sus evidencias existen y su sección de \texttt{DOCUMENTACION.md} está escrita.

% ---------- E0 ----------
\subsection{\texorpdfstring{\Ez\ }{E0 }Workspace y parámetros}
\ficha{cModelo}{Modelo}{Jueves 26, noche}{P1, P11}

\subsubsection{Qué hay que entender}
\begin{itemize}
  \item \textbf{Workspace colcon:} una carpeta con \texttt{src/}. \texttt{colcon build} compila cada paquete y deja el resultado en \texttt{install/}. \texttt{source install/setup.bash} le dice a la terminal dónde están.
  \item \textbf{Paquete \texttt{ament\_cmake} vs.\ \texttt{ament\_python}:}
  \begin{itemize}
    \item \texttt{ament\_cmake} sirve para archivos que solo se instalan (URDF, launch, configuración).
    \item \texttt{ament\_python} sirve para nodos en Python.
  \end{itemize}
  \item \textbf{Archivo único de parámetros:} todas las medidas y masas viven en \texttt{hexapod\_params.xacro}. El URDF las lee y la cinemática también (vía \texttt{robot\_description}). Cambiar un largo es cambiar una línea.
\end{itemize}

\subsubsection{Tareas}
\begin{enumerate}[label=T\arabic*]
  \item Crear los paquetes:
  \begin{itemize}
    \item \texttt{hexapod\_description}: el modelo;
    \item \texttt{hexapod\_control}: cinemática, marcha y teleoperación;
    \item \texttt{hexapod\_bringup}: Gazebo, controladores y launch.
  \end{itemize}
  \item Escribir \texttt{hexapod\_params.xacro} con las masas de la sección 4. Los largos parten de la cota $d \leq 8{,}8$ cm.
  \item Copiar desde Octavius las macros de inercia, los linters y \texttt{make\_release.sh}.
\end{enumerate}

\subsubsection{Entregables}
\begin{tabularx}{\linewidth}{@{}l Y Y@{}}
\toprule \textbf{} & \textbf{Entregable} & \textbf{Evidencia}\\ \midrule
D1 & 3 paquetes compilando & \texttt{colcon build} sin errores ni advertencias\\
D2 & Parámetros con origen de cada dato & Cada valor comentado: medido, hoja o hipótesis\\
\bottomrule
\end{tabularx}

% ---------- E1 ----------
\subsection{\texorpdfstring{\Eu\ }{E1 }Cinemática de una pata}
\ficha{cControl}{Control}{Jueves 26 noche -- viernes 27 mañana}{base de P8}

\subsubsection{Qué hay que entender}
Cada pata tiene 3 articulaciones:
\begin{itemize}
  \item \textbf{coxa} ($\theta_1$): gira en torno al eje vertical; mueve la pata adelante y atrás;
  \item \textbf{fémur} ($\theta_2$): sube y baja la pata;
  \item \textbf{tibia} ($\theta_3$): es la rodilla.
\end{itemize}
Los largos son $L_1$ (coxa), $L_2$ (fémur) y $L_3$ (tibia). El marco de la pata tiene $x$ hacia afuera del cuerpo y $z$ hacia arriba.

\textbf{Cinemática directa (FK).} Dados los ángulos, ¿dónde queda el pie?
\begin{align}
r &= L_1 + L_2\cos\theta_2 + L_3\cos(\theta_2+\theta_3)\\
x &= r\cos\theta_1,\qquad y = r\sin\theta_1,\qquad z = L_2\sin\theta_2 + L_3\sin(\theta_2+\theta_3)
\end{align}

\textbf{Cinemática inversa (IK).} Dado dónde quiero el pie, ¿qué ángulos necesito? Es lo que usa la marcha 50 veces por segundo.
\begin{align}
\theta_1 &= \operatorname{atan2}(y,\,x)\\
\rho &= \sqrt{x^2+y^2} - L_1,\qquad D^2 = \rho^2 + z^2\\
\cos\theta_3 &= \frac{D^2 - L_2^2 - L_3^2}{2L_2L_3},\qquad \theta_3 = -\arccos(\cdot)\quad\text{(rodilla hacia arriba)}\\
\theta_2 &= \operatorname{atan2}(z,\,\rho) - \operatorname{atan2}\!\big(L_3\sin\theta_3,\; L_2 + L_3\cos\theta_3\big)
\end{align}
La ecuación del codo es la \textbf{ley de cosenos} aplicada al triángulo fémur--tibia--recta al pie. Solo hay solución si $|L_2-L_3| \leq D \leq L_2+L_3$, es decir, si el pie está al alcance. Si no lo está, la función lo informa en vez de devolver ángulos inválidos.

\subsubsection{Tareas}
\begin{enumerate}[label=T\arabic*]
  \item \texttt{kinematics.py} con \texttt{fk()} e \texttt{ik()}, sin ROS, para poder probarlo solo.
  \item Test: recorrer miles de puntos alcanzables y verificar que FK(IK($p$)) $= p$.
  \item Macro \texttt{leg.xacro} y una pata sola en RViz con 3 sliders.
\end{enumerate}

\subsubsection{Entregables}
\begin{tabularx}{\linewidth}{@{}l Y Y@{}}
\toprule \textbf{} & \textbf{Entregable} & \textbf{Evidencia}\\ \midrule
D3 & FK/IK probadas & Test pasa con error $< 0{,}1$ mm\\
D4 & Pata en RViz coincide con FK & Captura: pose de sliders y posición del pie en TF iguales a FK\\
\bottomrule
\end{tabularx}

\textbf{Para la presentación:} el dibujo de la pata con sus 3 ángulos, las ecuaciones de la IK y el resultado del test.

% ---------- E2 ----------
\subsection{\texorpdfstring{\Ed\ }{E2 }Modelo completo en RViz}
\ficha{cModelo}{Modelo}{Viernes 27, mañana y tarde}{P2, P3, P4, P5}

\subsubsection{Qué hay que entender}
\begin{itemize}
  \item \textbf{URDF:} describe al robot como un árbol de eslabones (\emph{links}) unidos por articulaciones (\emph{joints}). Cada eslabón tiene tres partes:
  \begin{itemize}
    \item \texttt{visual}: lo que se ve;
    \item \texttt{collision}: lo que choca;
    \item \texttt{inertial}: masa e inercia.
  \end{itemize}
  \item \textbf{XACRO:} URDF con variables y macros. La macro de pata se llama 6 veces, cambiando el nombre y el ángulo de montaje $\varphi_i \in \{\pm30^\circ, \pm90^\circ, \pm150^\circ\}$.
  \item \textbf{Inercia de una caja} de masa $m$ y lados $a, b, c$. Gazebo la necesita realista; si no, el robot tiembla o “explota”:
  \begin{equation}
  I_{xx} = \tfrac{m}{12}(b^2+c^2),\quad I_{yy} = \tfrac{m}{12}(a^2+c^2),\quad I_{zz} = \tfrac{m}{12}(a^2+b^2)
  \end{equation}
  \item \textbf{Límites del servo en el URDF:}
  \begin{itemize}
    \item velocidad: $60^\circ / 0{,}14\,\text{s} \approx 7{,}5$ rad/s;
    \item esfuerzo: 1,08 N·m;
    \item rango de cada articulación: el mecánico.
  \end{itemize}
  \item \textbf{\texttt{robot\_state\_publisher}:} toma el URDF y \texttt{/joint\_states}, y publica la TF, que es la posición de cada eslabón. \textbf{\texttt{joint\_state\_publisher\_gui}:} publica \texttt{/joint\_states} desde sliders.
\end{itemize}

\subsubsection{Tareas}
\begin{enumerate}[label=T\arabic*]
  \item \texttt{body.xacro}: hexágono regular con radio $R$ del archivo de parámetros.
  \item \texttt{components.xacro}: Raspberry, 2 PCA9685 y batería (abajo, para bajar el centro de masa) como cajas con su masa.
  \item Cada servo como caja de 55 g en su posición real dentro de la pata.
  \item \texttt{sensors.xacro}: \texttt{imu\_link} al centro y \texttt{camera\_link} + \texttt{camera\_optical\_link} en la cara frontal.
  \item \texttt{display.launch.py} y \texttt{display.rviz} con RobotModel y TF.
\end{enumerate}

\subsubsection{Entregables}
\begin{tabularx}{\linewidth}{@{}l Y Y@{}}
\toprule \textbf{} & \textbf{Entregable} & \textbf{Evidencia}\\ \midrule
D5 & URDF válido con 18 articulaciones & \texttt{check\_urdf} OK y árbol en PDF\\
D6 & RViz con 18 sliders & Captura moviendo una pata\\
D7 & Masa total del modelo $\approx$ masa estimada & Suma de masas del URDF $= 1{,}87 \pm 0{,}05$ kg\\
\bottomrule
\end{tabularx}

% ---------- E3 ----------
\subsection{\texorpdfstring{\Et\ }{E3 }De pie en Gazebo}
\ficha{cSim}{Simulación}{Viernes 27, tarde y noche --- \textbf{etapa de mayor riesgo}}{P6, base de P8}

\subsubsection{Qué hay que entender}
\begin{itemize}
  \item \textbf{Base flotante:} a diferencia de Octavius, nada fija el cuerpo al mundo. El robot solo se sostiene si:
  \begin{itemize}
    \item los pies tocan el suelo;
    \item la fricción impide que resbalen;
    \item los servos aguantan el peso.
  \end{itemize}
  \item \textbf{Contacto:} el pie es una \textbf{esfera}, porque toca el suelo en un punto y es estable. Se le da fricción alta ($\mu \approx 1$) y un contacto algo blando para que no rebote.
  \item \textbf{\texttt{gz\_ros2\_control}:} conecta \texttt{ros2\_control} con el motor de física. Para cada articulación expone:
  \begin{itemize}
    \item una interfaz de \textbf{comando} de posición (lo que pedimos);
    \item interfaces de \textbf{estado} de posición y velocidad (lo medido, como un encoder ideal).
  \end{itemize}
  \item \textbf{Postura inicial suave:} el robot aparece con las patas recogidas. Las patas bajan de a poco hasta levantar el cuerpo a la altura de trabajo. Si se pidiera la postura de golpe, el robot saltaría.
\end{itemize}

\subsubsection{Tareas}
\begin{enumerate}[label=T\arabic*]
  \item \texttt{ros2\_control.xacro} y \texttt{controllers.yaml} con \texttt{joint\_state\_broadcaster} y \texttt{legs\_controller} (18 articulaciones).
  \item \texttt{gazebo.launch.py}: mundo plano, \emph{spawn} sobre el suelo y carga de controladores.
  \item Nodo o rutina de postura: interpolar desde las patas recogidas hasta la postura de pie en 2--3 s.
  \item Ajustar fricción, contacto y paso de física hasta que se sostenga.
\end{enumerate}

\subsubsection{Entregables}
\begin{tabularx}{\linewidth}{@{}l Y Y@{}}
\toprule \textbf{} & \textbf{Entregable} & \textbf{Evidencia}\\ \midrule
D8 & Robot de pie en Gazebo & 60 s sin hundirse; altura del cuerpo estable ($\pm 2$ mm) en \texttt{/tf}\\
D9 & Controladores activos & \texttt{ros2 control list\_controllers}: ambos \emph{active}\\
\bottomrule
\end{tabularx}

\textbf{Puntos por tratar.} Si el robot se hunde o tiembla, revisar en este orden:
\begin{enumerate}
  \item las inercias;
  \item las colisiones duplicadas entre eslabones vecinos;
  \item la fricción;
  \item el paso de física.
\end{enumerate}
Si sigue fallando el viernes en la noche, se reduce la masa del cuerpo en el modelo solo para validar el resto, y se anota en la documentación.

% ---------- E4 ----------
\subsection{\texorpdfstring{\Ec\ }{E4 }Marcha trípode y teclado}
\ficha{cControl}{Control}{Sábado 28, mañana y tarde}{P8}

\subsubsection{Qué hay que entender}
\textbf{Marcha trípode.} Las patas se dividen en dos grupos que se alternan:
\begin{itemize}
  \item grupo A: LF, RM, LR;
  \item grupo B: RF, LM, RR.
\end{itemize}
Mientras un grupo apoya, el otro vuela. Siempre hay un triángulo de 3 patas en el suelo.

Cada pata $i$ tiene una fase $\phi_i \in [0,1)$ que avanza con el tiempo, con periodo $T$. El grupo B va desfasado medio ciclo. Con fracción de apoyo $\beta = 0{,}5$:
\begin{align}
\text{apoyo } (\phi < \beta):&\quad \mathbf p_i = \mathbf p_{0,i} + \Big(\tfrac12 - \tfrac{\phi}{\beta}\Big)\boldsymbol\Delta_i\\
\text{vuelo } (\phi \geq \beta):&\quad u = \tfrac{\phi-\beta}{1-\beta},\quad \mathbf p_i = \mathbf p_{0,i} + \Big(u-\tfrac12\Big)\boldsymbol\Delta_i + h\sin(\pi u)\,\hat z
\end{align}
\begin{itemize}
  \item En el \textbf{apoyo}, el pie se mueve hacia atrás respecto del cuerpo, y eso empuja al robot hacia adelante.
  \item En el \textbf{vuelo}, el pie vuelve adelante por un arco de altura $h$.
\end{itemize}

El \textbf{vector de paso} sale de la velocidad pedida. Incluye el giro, porque al girar cada pata recorre un arco distinto según su posición $\mathbf p_{0,i}$:
\begin{equation}
\boldsymbol\Delta_i = \beta T\,\big(\mathbf v + \omega_z\,\hat z \times \mathbf p_{0,i}\big)
\end{equation}
Cada $\mathbf p_i$ (en el marco del cuerpo) se pasa al marco de su pata. Después la IK de \Eu\ entrega los 3 ángulos.

\subsubsection{Tareas}
\begin{enumerate}[label=T\arabic*]
  \item \texttt{gait.py}, sin ROS: fase, trayectorias y vector de paso, con tests.
  \item \texttt{gait\_node.py}: escucha \texttt{/hexapod/cmd\_vel} y publica 18 ángulos a 50 Hz en \texttt{/legs\_controller/commands}.
  \item \texttt{keyboard\_teleop.py}, adaptado de Octavius: las teclas ajustan $v_x$, $v_y$ y $\omega_z$.
\end{enumerate}

\subsubsection{Entregables}
\begin{tabularx}{\linewidth}{@{}l Y Y@{}}
\toprule \textbf{} & \textbf{Entregable} & \textbf{Evidencia}\\ \midrule
D10 & Camina con teclado & Video: avanza, retrocede y gira\\
D11 & Velocidad real $\approx$ pedida & Desplazamiento medido en \texttt{/tf} durante 10 s, con error menor al 20\,\%\\
\bottomrule
\end{tabularx}

\textbf{Para la presentación:} el diagrama de los dos grupos de patas en el tiempo y el video caminando.

% ---------- E5 ----------
\subsection{\texorpdfstring{\Eci\ }{E5 }Sensores: IMU y cámara}
\ficha{cSens}{Sensores}{Sábado 28, tarde}{P5, P7}

\subsubsection{Qué hay que entender}
\begin{itemize}
  \item \textbf{IMU} (\texttt{sensor\_msgs/Imu}): entrega la orientación del cuerpo (cuaternión), la velocidad angular y la aceleración. En un hexápodo sirve para saber si el cuerpo está inclinado. El plugin de Gazebo acepta ruido, que se configura con valores de la hoja del BNO085.
  \item \textbf{Cámara} (\texttt{sensor\_msgs/Image}): se modela como cámara estenopeica (\emph{pinhole}). Con un ancho de imagen $W$ y un campo de visión horizontal (HFOV), la distancia focal en píxeles es:
  \begin{equation}
  f_x = \frac{W/2}{\tan(\text{HFOV}/2)} = \frac{320}{\tan 33^\circ} \approx 493\ \text{px}\quad (W = 640,\ \text{HFOV} = 66^\circ)
  \end{equation}
  La resolución se reduce respecto a la cámara real (12 MP) para no cargar el PC.
  \item \textbf{\texttt{ros\_gz\_bridge}:} traduce los tópicos de Gazebo a ROS\,2. Se configura en \texttt{gz\_bridge.yaml}.
  \item \textbf{\texttt{camera\_optical\_link}:} la convención REP-103 pide que la cámara mire hacia $+z$ con $x$ a la derecha, distinto del resto del robot. Por eso se agrega un marco extra.
\end{itemize}

\subsubsection{Entregables}
\begin{tabularx}{\linewidth}{@{}l Y Y@{}}
\toprule \textbf{} & \textbf{Entregable} & \textbf{Evidencia}\\ \midrule
D12 & IMU publicando & \texttt{ros2 topic hz /hexapod/imu} $\approx$ 100 Hz; al inclinar el robot cambia la orientación\\
D13 & Cámara publicando & \texttt{ros2 topic hz /hexapod/camera/image} $\approx$ 15 Hz; imagen visible en RViz\\
\bottomrule
\end{tabularx}

% ---------- E6 ----------
\subsection{\texorpdfstring{\Es\ }{E6 }micro-ROS: joystick ESP32}
\ficha{cMicro}{micro-ROS}{Sábado 28, noche --- \textbf{la prueba con hardware queda para cuando se compre la ESP32}}{P9}

\subsubsection{Qué hay que entender}
\begin{itemize}
  \item \textbf{micro-ROS:} una versión de ROS\,2 para microcontroladores. La ESP32 no habla DDS directamente. Habla con el \textbf{agente} (\texttt{micro\_ros\_agent}) por USB serial, y el agente la presenta al resto de ROS\,2 como un nodo más.
  \item \textbf{Qué hace la ESP32:} lee los dos ejes del joystick con el ADC, los normaliza a $[-1, 1]$ y los publica como \texttt{sensor\_msgs/Joy} en \texttt{/hexapod/joystick} a 20 Hz. Reutiliza el firmware de Octavius.
  \item \textbf{\texttt{joy\_to\_twist}:} convierte el joystick en velocidad del cuerpo, con una zona muerta $\delta$ para que el robot no se mueva con el joystick suelto:
  \begin{equation}
  \operatorname{dz}(a) = \begin{cases}0 & |a| < \delta\\ \operatorname{sign}(a)\,\dfrac{|a|-\delta}{1-\delta} & |a| \geq \delta\end{cases}\qquad v_x = v_{\max}\operatorname{dz}(a_y),\quad \omega_z = \omega_{\max}\operatorname{dz}(a_x)
  \end{equation}
\end{itemize}

\subsubsection{Entregables}
\begin{tabularx}{\linewidth}{@{}l Y Y@{}}
\toprule \textbf{} & \textbf{Entregable} & \textbf{Evidencia}\\ \midrule
D14 & \texttt{joy\_to\_twist} con test & Test de zona muerta y escalas pasa\\
D15 & Robot se mueve con el tópico a mano & \texttt{ros2 topic pub /hexapod/joystick ...} hace caminar al robot\\
D16 & Firmware compilado & Compila sin advertencias (Arduino CLI)\\
D17 & \emph{Con la ESP32:} joystick real mueve al robot & Video: nodo de la ESP32 en \texttt{ros2 node list} y robot caminando\\
\bottomrule
\end{tabularx}

% ---------- E7 ----------
\subsection{\texorpdfstring{\Esi\ }{E7 }Integración y pruebas}
\ficha{cInteg}{Integración}{Sábado 28, noche}{P10, P11}

\subsubsection{Tareas}
\begin{enumerate}[label=T\arabic*]
  \item \texttt{bringup.launch.py}, que levanta todo. Argumentos: \texttt{use\_rviz}, \texttt{use\_microros}, \texttt{serial\_port}, \texttt{world}.
  \item Test de consistencia: el URDF, \texttt{controllers.yaml} y \texttt{gz\_bridge.yaml} usan los mismos nombres de articulaciones y tópicos. Más los linters de ROS\,2.
  \item Prueba en limpio: instalar desde cero en un contenedor \texttt{ros:jazzy} siguiendo solo el README. \emph{Se recorta si falta tiempo.}
\end{enumerate}

\subsubsection{Entregables}
\begin{tabularx}{\linewidth}{@{}l Y Y@{}}
\toprule \textbf{} & \textbf{Entregable} & \textbf{Evidencia}\\ \midrule
D18 & Un comando levanta todo & \texttt{ros2 launch hexapod\_bringup bringup.launch.py}\\
D19 & Tests pasan & \texttt{colcon test-result}: 0 fallas\\
\bottomrule
\end{tabularx}

% ---------- E8 ----------
\subsection{\texorpdfstring{\Eo\ }{E8 }Estudio de dimensiones}
\ficha{cEstudio}{Estudio}{Domingo 29, mañana y tarde}{aporte propio (no es de la pauta)}

\subsubsection{Pregunta}
Con 18 MG996R, una LiPo 3S y la electrónica dada, ¿qué largos $L_1, L_2, L_3$ y qué radio $R$ permiten caminar con margen de torque, sin choques y de forma estable?

\subsubsection{Restricciones}
\begin{tabularx}{\linewidth}{@{}l Y@{}}
\toprule \textbf{Criterio} & \textbf{Condición}\\ \midrule
Torque del fémur & $F\,\rho \leq \tau_d$, donde $\rho$ es la distancia horizontal del eje del fémur al pie (incluye el peso propio de la pata)\\
Torque de la tibia & $F\,|L_3\cos(\theta_2+\theta_3)| \leq \tau_d$\\
Alcance & El paso completo $\pm\Delta/2$ cabe en el espacio de trabajo, dentro de los límites articulares\\
Altura libre & Altura del cuerpo $\geq h_{\min}$ (propuesta: 6 cm)\\
Sin choques & Patas vecinas, a $60^\circ$ entre sí, no se tocan en ningún punto del ciclo\\
Estabilidad & Distancia del centro de masa a los bordes del triángulo de apoyo $\geq S_{\min}$\\
Layout & Batería, Raspberry y PCA9685 caben dentro del hexágono, con el centro de masa al centro\\
\bottomrule
\end{tabularx}

\subsubsection{Tareas}
\begin{enumerate}[label=T\arabic*]
  \item \texttt{analysis/dimension\_study.py}: recorre combinaciones de $(L_1, L_2, L_3, R)$, aplica las restricciones y guarda las que cumplen.
  \item Gráficos:
  \begin{itemize}
    \item mapa de torque del fémur según $L_2$ y $L_3$;
    \item espacio de trabajo de una pata;
    \item margen de estabilidad.
  \end{itemize}
  \item Elegir la mejor combinación, por mayor margen mínimo, y copiarla a \texttt{hexapod\_params.xacro}.
  \item \emph{(Bajo la línea de corte.)} Verificar en Gazebo.
  \begin{itemize}
    \item \textbf{Primero, una prueba de confianza:} colgar una pata con una carga conocida y comparar el esfuerzo que informa \texttt{gz\_ros2\_control} con $m\,g\,d$.
    \item Si coincide, registrar los torques durante la marcha y compararlos con $\tau_d$.
    \item Si no coincide, quedarse con el valor analítico y decirlo.
  \end{itemize}
  \item \texttt{docs/especificacion\_diseno.md}: cotas para maquinar (de eje a eje), posición de cada componente y márgenes obtenidos.
\end{enumerate}

\subsubsection{Entregables}
\begin{tabularx}{\linewidth}{@{}l Y Y@{}}
\toprule \textbf{} & \textbf{Entregable} & \textbf{Evidencia}\\ \midrule
D20 & Cotas ideales & Tabla con $L_1, L_2, L_3, R$ y el margen de cada restricción\\
D21 & Gráficos del estudio & 3 figuras en \texttt{docs/img/}\\
D22 & Especificación de diseño & \texttt{especificacion\_diseno.md} con cotas y layout\\
\bottomrule
\end{tabularx}

% ---------- E9 ----------
\subsection{\texorpdfstring{\En\ }{E9 }Documentación y entrega}\label{sec:entrega}
\ficha{cEntrega}{Entrega}{Cada cierre de etapa + domingo 29, tarde y noche}{P12, P13, P14}

\subsubsection{Tareas}
\begin{enumerate}[label=T\arabic*]
  \item \texttt{DOCUMENTACION.md}: una sección por etapa, escrita al cerrarla. Para cada componente: qué es, por qué se usa, qué datos entran y salen, y cómo se procesan.
  \item Tutorial: \texttt{DOCUMENTACION.md} $\rightarrow$ \texttt{.docx} y \texttt{.pdf}.
  \item Presentación Beamer con la plantilla USACH de Octavius ($\leq$ 15 min): \texttt{.tex} editable + \texttt{.pdf}.
  \item \texttt{make\_release.sh} $\rightarrow$ \texttt{cesar\_olivares\_2026\_2.zip}, sin \texttt{build/}, \texttt{install/} ni \texttt{log/}.
\end{enumerate}

\subsubsection{Checklist de formato de la presentación}
\begin{itemize}
  \item Membrete y logo USACH.
  \item Título, estudiante, profesor (José Pascal) y fecha.
  \item Numeración de páginas, tablas, figuras y ecuaciones.
  \item Letra uniforme y figuras con buena resolución.
  \item Índice, introducción y conclusiones.
  \item Referencias IEEE.
\end{itemize}

\subsubsection{Guion sugerido (15 min)}
\begin{tabularx}{\linewidth}{@{}l r Y@{}}
\toprule \textbf{Bloque} & \textbf{Min} & \textbf{Contenido}\\ \midrule
Introducción & 1,5 & Qué es un hexápodo, objetivo y alcance\\
Modelo & 3 & Cuerpo hexagonal, macro de pata, componentes reales, RViz\\
Cinemática y marcha & 3,5 & FK/IK, trípode, video caminando\\
Sensores y micro-ROS & 2,5 & IMU, cámara, flujo del joystick\\
Estudio de dimensiones & 3 & Restricciones, gráficos, cotas finales\\
Conclusiones & 1,5 & Qué se logró, qué falta (ESP32, robot real)\\
\bottomrule
\end{tabularx}

% ====================== 8. CONTROL ======================
\Needspace{14\baselineskip}
\section{Control de avance}

\textbf{Avance} $=$ entregables cumplidos $/$ 22. Un entregable cuenta solo si su evidencia existe. D17 depende de la ESP32 y queda fuera del total de esta semana, así que el máximo esta semana es 21/22.

\textbf{Revisión al final de cada día} (en vez de mensual, por el plazo), con semáforo por etapa:

\begin{tabularx}{\linewidth}{@{}l Y@{}}
\toprule
\pill{sVerde}{Verde} & Todo lo del día cumplido.\\
\pill{sAmbar}{Ámbar} & Atraso de un bloque o menos. Se mueve al bloque siguiente.\\
\pill{sRojo}{Rojo} & Falta la mitad o más de la etapa, o peligra un punto de la pauta. Se aplica la línea de corte.\\
\bottomrule
\end{tabularx}

\textbf{Cierre de etapa.} Una nota corta en \texttt{DOCUMENTACION.md} con:
\begin{itemize}
  \item qué se hizo;
  \item las evidencias;
  \item los problemas encontrados y cómo se resolvieron;
  \item lo que hay que estudiar para la presentación.
\end{itemize}

% ====================== 9. PUNTOS ======================
\Needspace{14\baselineskip}
\section{Puntos por tratar}

\begin{tabularx}{\linewidth}{@{}Y l Y@{}}
\toprule \textbf{Punto} & \textbf{Afecta} & \textbf{Qué se hace mientras}\\ \midrule
Fecha y hora exactas de entrega & Cronograma & Se asume domingo 29 nov, 23:59\\
ESP32 + joystick por comprar & P9 & Código listo; prueba con el tópico a mano\\
¿Los 55 g del servo incluyen los adaptadores coaxiales? & Masas & Se asume que sí\\
Medidas L × A × H de la batería & Layout (\Eo) & Se usa un tamaño típico de 3S 5200 mAh, marcado como hipótesis\\
Masa de la estructura impresa & Torque (\Eo) & Hipótesis de 400 g, que se recalcula con las cotas finales\\
\bottomrule
\end{tabularx}

% ====================== 10. GLOSARIO ======================
\Needspace{14\baselineskip}
\section{Glosario}

\begin{tabularx}{\linewidth}{@{}l Y@{}}
\toprule
Coxa, fémur, tibia & Los tres segmentos de una pata, de la base hacia el pie.\\
FK / IK & Cinemática directa (de ángulos a posición del pie) e inversa (de posición a ángulos).\\
Marcha trípode & Marcha con dos grupos de 3 patas que se alternan entre apoyo y vuelo.\\
Base flotante & El cuerpo no está fijo al mundo; se sostiene por el contacto de las patas.\\
URDF / XACRO & Formato que describe el robot. XACRO agrega variables y macros.\\
TF & Árbol de marcos de referencia con la posición de cada eslabón en el tiempo.\\
\texttt{ros2\_control} & Marco de ROS\,2 que conecta controladores con el hardware (o con Gazebo).\\
micro-ROS & ROS\,2 para microcontroladores; se comunica a través de un agente en el PC.\\
Torque de bloqueo & Torque máximo del servo trabado. Se diseña con la mitad.\\
Margen de estabilidad & Distancia del centro de masa a los bordes del triángulo de apoyo.\\
\bottomrule
\end{tabularx}

\Needspace{14\baselineskip}
\section{Enlaces útiles}
\begin{itemize}
  \item ROS\,2 Jazzy: \url{https://docs.ros.org/en/jazzy/}
  \item ros2\_control: \url{https://control.ros.org/jazzy/}
  \item Gazebo Harmonic: \url{https://gazebosim.org/docs/harmonic/}
  \item micro-ROS: \url{https://micro.ros.org/}
\end{itemize}

\end{document}
